\section{Discussion}
\label{sec:discussion-ra}

The reversible-action discipline extends the amendment-side type-level
rule of the prior substrate \cite{programmableSovereigntyAnonymous} to
the response side, for the class of actions whose effects are
reversible by an inverse executor. The interpretation of the extension
is constrained by what the substrate proves and what it does not. Three
threads of discussion follow: the relationship to constitutional
emergency powers, the relationship to agentic-AI oversight, and the
relationship to the substrate's notion of time.

\paragraph{Emergency powers and bounded action.}
A natural framing for time-bounded constitutional change is the
literature on emergency powers, particularly the post-9/11 work on
deliberative versus emergency authority. The framing is partially
apt: a TTL-bounded amendment looks like an emergency exception, and
the auto-reversion at TTL expiry looks like a sunset clause. The
framing is also partially misleading: real emergency powers are
notable for the absence of reliable rollback. A government that
invokes emergency authority does not in practice undo the actions
taken during the emergency; the very property the TTL discipline
formalizes is the property emergency-powers regimes have struggled to
guarantee. The substrate therefore does not claim to formalize
emergency powers; it formalizes a stricter discipline (timed reversal
with type-level invariant) that lies inside the space emergency-
powers scholarship surveys.

The deeper point is that the discipline is conservative. The prior
substrate's amendment refinement obligation
\cite{programmableSovereigntyAnonymous} already requires backward
refinement: an amendment cannot widen admission on receipts already
admitted. The TTL-bounded extension keeps the refinement obligation per
step and adds the reversion. A constitutional change that genuinely
cannot be reverted (debt restructuring, retroactive statute, founder
override, treaty repudiation) does not fit the TTL-bounded class; the
prior substrate records such events as crisis artifacts that break
refinement and are signed as such. The TTL-bounded class is a strict
subclass of the prior substrate's amendment space.

\paragraph{Agentic-AI oversight.}
An autonomous agent taking a destructive action without a paired
rollback is exactly the bounded-action failure mode the substrate
forbids at the type level. A typed substrate that refuses to
construct an executive action without a TTL witness and an action
class tag is a primitive other safety systems can cite without
adopting the substrate as a whole. Constitutional AI conditions
outputs at training time; capability evaluations bound what a model
is likely to do in advance; the substrate here gates whether a
specific cross-boundary action enters the receipt log without the
counterparty's signature.

The bilateral-cosignature reduction on the destructive subclass is
the natural primitive for human-in-the-loop deployment of agentic
systems: a destructive action authored by an agent admits only when
the deploying organization's polity also signs. The reduction does
not solve the operational-discipline problem already named by the
prior substrate \cite{programmableSovereigntyAnonymous}
(a single actor operating both keys collapses the bilateral envelope
to one-of-one), but it gives the discipline a typed obligation rather
than a policy file. The Anthropic alignment-evaluation work on
alignment-faking and situational-awareness is the upstream
motivation; the reversible-action substrate is the downstream
discipline.

\paragraph{Time and the substrate.}
The substrate treats time as a monotonic integer counter. The TTL
invariant is a duration in counter units, and the auto-reversion at
TTL expiry is a substrate-level event triggered by the counter
advancing past the expiry instant. The substrate model abstracts the
real-time scheduler that advances the counter; the scheduler is a
deployment-side obligation. The implication for the threat model is
that a compromised counter falsifies the TTL claim in the same way a
compromised signing key falsifies admission. The substrate does not
trust the endpoint's wall clock; it trusts the counter's monotonicity
and the scheduler's faithfulness in advancing it.

The substrate's notion of partition (the partition-contingency mode on
the trust ladder, carried by the prior substrate
\cite{programmableSovereigntyAnonymous}) introduces a second
relationship to time. A partition is a window during which the
substrate cannot exchange receipts with peer polities; the substrate
records the partition's start and end as receipts and conditions
admission during the partition on the ladder's partition-contingency
rules. A TTL-bounded amendment whose TTL spans a partition is a
substrate-level case the model does not yet cover: the auto-reversion
at TTL expiry should fire even if the substrate is partitioned, but
the reversion's receipt cannot be exchanged with peers until the
partition ends. The substrate model treats this as a deployment-side
obligation rather than as a substrate-level theorem.

\paragraph{What the substrate does not extend.}
The reversible-action discipline does not extend the type-level rule
to actions whose effects are not reversible by an inverse executor.
Process-tree terminate, network isolate, and grant revoke are
destructive: the substrate's response to these is to require bilateral
cosignature at admission, not to record a rollback. The substrate makes
the distinction between reversible and destructive explicit at the type
level, but it does not turn a destructive action into a reversible one.
The category boundary is empirical (existence of an inverse executor)
rather than constitutional.
